\part{Group Theory}
\chapter{Discrete Symmetry Groups}
\section{Finite subgroups of the orthogonal groups}
\begin{importanttheorem}{}{}
	The finite subgroups of $SO(3)$, up to isomorphism, are:
	\begin{enumerate}
		\item The cyclic groups $C_n$,
		\item The dihedral groups $D_n$,
		\item The tetrahedral group $T \cong A_4$,
		\item The cube group $W \cong S_4$,
		\item The icosahedral group $I \cong A_5$.
	\end{enumerate}
\end{importanttheorem}
\begin{importantcorollary}{}{}
	The finite subgroups of $O(3)$, up to isomorphism, are:
	\begin{enumerate}
		\item The finite subgroups of $SO(3)$:
		\begin{enumerate}
			\item The cyclic groups $C_n$,
			\item The dihedral groups $D_n$,
			\item The tetrahedral group $T \cong A_4$,
			\item The cube group $W \cong S_4$,
			\item The icosahedral group $I \cong A_5$.
		\end{enumerate}
		\item Products of these groups with $C_2$.
	\end{enumerate}
\end{importantcorollary}
\section{Lattice Groups}
\begin{definition}
	An $n$-dimensional real \tib{lattice} is a subset of $\bR^n$ of the form
	\begin{align*}
		\Lambda = \lr{\sum_{i = 1}^n k_i.v_i \mid k_i \in \bZ},
	\end{align*}
	where $v_1, \hdots, v_n$ is a basis of $\bR^n$. We say that the lattice $\Lambda$ is \tib{generated by the matrix} $M = (v_1, \hdots, v_n)$.
\end{definition}
\begin{importanttheorem}{Crystallographic restriction theorem}{}
	Every rotational symmetry of a two- or three-dimensional lattice has order $1$, $2$, $3$, $4$, or $6$.
\end{importanttheorem}
\section{The Two-Dimensional Point Groups}
\begin{definition}
	An $n$-dimensional \tib{crystallographic point group} is a subgroup of $O(n)$ mapping an $n$-dimensional lattice to itself while leaving a given point unchanged.
\end{definition}
\begin{importanttheorem}{}{}
	The two-dimensional crystallographic point groups are:
	\begin{enumerate}
		\item The cyclic groups $C_1, C_2, C_3, C_4$ and $C_6$,
		\item The dihedral groups $D_1, D_2, D_3, D_4$ and $D_6$.
	\end{enumerate}
	We have $D_1 \cong C_2$, so there are ten point groups up to equality and nine point groups up to equality.
\end{importanttheorem}
\begin{proofsketch}
	Follows immediately from the classification of finite subgroups of $O(2)$ and the crystallographic restriction theorem.
\end{proofsketch}
\iffalse
\section{Lattice Systems}
\begin{importanttheorem}{Two-dimensional Bravais lattices}{}
	There are five lattice types in two dimensions:
	\begin{enumerate}
		\item The \tib{oblique} or \tib{scalene triangular} lattice,
		\item The \tib{rhombic}, \tib{centered rectangular} or \tib{isosceles triangular} lattice,
		\item The \tib{rectangular}, \tib{centered rhombic} or \tib{right triangular} lattice,
		\item The \tib{square} or \tib{right isosceles triangular} lattice,
		\item The \tib{hexagonal} or \tib{equilateral triangular} lattice.
	\end{enumerate}
\end{importanttheorem}
Before we get to the classification of three-dimensional Bravais lattices, I need to establish some vocabulary - three-dimensional lattices are named based on the smallest repeating unit within the lattice that still captures its full symmetry group. The possibilities here turn out to be:
\begin{enumerate}
	\item \tib{triclinic systems} ("three skew angle systems"), given by repeating parallelepipeds where all side lengths are unequal, all angles are unequal, and none of the angles are right angles.
	\item \tib{monoclinic systems} ("one skew angle systems"), given by repeating parallelepipeds where all side lengths are unequal, two angles are right angles, and the third isn't. "Biclinic systems", i.e. systems with exactly one right angle, are ignored within our classification, since they don't add additional symmetries beyond those of a triclinic systems.
	\item \tib{orthorhombic systems}, given by repeating cuboids with three unequal side lengths,
	\item \tib{tetragonal systems}, given by repeating cuboids with square bases,
	\item \tib{hexagonal systems}, given by repeating hexagonal prisms,
	\item \tib{rhombohedral systems}, given by repeating rhombohedrons (i.e. parallelepipeds with three equal side lengths, but three different angles)
	\item and \tib{cubic systems}, given by repeating cubes.
\end{enumerate}
\fi
\section{The Three-Dimensional Point Groups}
\begin{importanttheorem}{}{}
	There are eighteen three-dimensional crystallographic point groups up to isomorphism, and they are:
	\begin{enumerate}
		\item Five cyclic groups:
		\begin{itemize}
			\item $C_1$
			\item $C_2 \cong D_1$
			\item $C_3$
			\item $C_4$
			\item $C_6 \cong C_3 \times C_2$
		\end{itemize}
		\item Four non-cyclic dihedral groups:
		\begin{itemize}
			\item $D_2 \cong C_2 \times C_2$
			\item $D_3$
			\item $D_4$ 
			\item $D_6 \cong D_3 \times C_2$
		\end{itemize}
		\item Two remaining products of cyclic groups with $C_2$: 
		\begin{itemize}
			\item $C_4 \times C_2$
			\item $C_6 \times C_2$
		\end{itemize}
		\item Three remaining products of dihedral groups with $C_2$: 
		\begin{itemize}
			\item $D_2 \times C_2$
			\item $D_4 \times C_2$ 
			\item $D_6 \times C_2$
		\end{itemize}
		\item Four remaining groups consisting of permutation groups and their products with $C_2$:
		\begin{itemize}
			\item $A_4$
			\item $A_4 \times C_2$
			\item $S_4$
			\item $S_4 \times C_2$
		\end{itemize}
	\end{enumerate}
\end{importanttheorem}
\begin{importanttheorem}{}{}
	The following table gives a full classification of three-dimensional crystallographic point groups up to identity:
\end{importanttheorem}
\begin{center}
	\begin{tabular}{c|c|cc|c}
		Crystal System & Crystal Class & Example Figure & HM Notation & Isomorphism\\\hline
		\multirow{2}{*}{triclinic} & triclinic-pedial & & $1$ & $C_1$\\\cline{3-5}
		& triclinic-pinacoidal & & $\ol 1$ & \multirow{3}{*}{$C_2$}\\\cline{1-4}
		\multirow{3}{*}{monoclinic} & monoclinic-sphenoidal & & $2$ \\\cline{3-4}
		& monoclinic-domatic & & $m$\\\cline{3-5}
		& monoclinic-prismatic & & $2/m$ & \multirow{3}{*}{$D_2$}\\\cline{1-4}
		\multirow{3}{*}{orthorhombic} & orthorhombic-disphenoidal & & $222$ &\\\cline{3-4}
		& orthorhombic-pyramidal & & $mm2$ &\\\cline{3-5}
		& orthorhombic-dipyramidical & & $mmm$ & $D_2 \times C_2$\\\cline{1-5}
		\multirow{7}{*}{tetragonal} & tetragonal-pyramidal & & $4$ & \multirow{2}{*}{$C_4$}\\\cline{3-4}
		& tetragonal-disphenoidal & & $\ol 4$ &\\\cline{3-5}
		& tetragonal-dipyramidal & & $4/m$ & $C_4 \times C_2$\\\cline{3-5}
		& tetragonal-trapezoidal & & $422$ & \multirow{3}{*}{$D_4$}\\\cline{3-4}
		& ditetragonal-pyramidal & & $4mm$ &\\\cline{3-4}
		& tetragonal-scalenoidal & & $\ol 42m$ &\\\cline{3-5}
		& ditetragonal-dipyramidal & & $4/mmm$ & $D_4 \times C_2$\\\cline{1-5}
		\multirow{5}{*}{trigonal} & trigonal-pyramidal & & $3$ & $C_3$\\\cline{3-5}
		& rhombohedral & & $\ol 3$ & $C_6$\\\cline{3-5}
		& trigonal-trapezoidal & & $32$ & \multirow{2}{*}{$D_3$}\\\cline{3-4}
		& ditrigonal-pyramidal & & $3m$ &\\\cline{3-5}
		& ditrigonal-scalenohedral & & $\ol 3m$ & $D_6$\\\cline{1-5}
		\multirow{7}{*}{hexagonal} & hexagonal-pyramidal & & $6$ & \multirow{2}{*}{$C_6$}\\\cline{3-4}
		& trigonal-dipyramidal & & $\ol 6$ &\\\cline{3-5}
		& hexagonal-dipyramidal & &$6/m$ & $C_6 \times C_2$\\\cline{3-5}
		& hexagonal-trapezohedral & & $622$ & \multirow{3}{*}{$D_6$}\\\cline{3-4}
		& dihexagonal-pyramidal & & $6mm$ &\\\cline{3-4}
		& ditrigonal-dipyramidal & & $\ol 6 m 2$\\\cline{3-5}
		& dihexagonal-dipyramidal & & $6/mmm$ & $D_6 \times C_2$\\\cline{1-5}
		\multirow{5}{*}{cubic} & tetartoidal & & $23$ & $A_4$\\\cline{3-5}
		& diploidal & & $m\ol 3$ & $A_4 \times C_2$\\\cline{3-5}
		& gyroidal & & $432$ & \multirow{2}{*}{$S_4$}\\\cline{3-4}
		& hextetrahedral & & $\ol 4 3 m$ &\\\cline{3-5}
		& hexoctahedral & & $m \ol 3 m$ & $S_4 \times C_2$\\
	\end{tabular}
\end{center}
\section{Classification of Lattices}
\section{Wallpaper Groups}
\section{Three-Dimensional Space Groups}